\documentclass[fleqn,10pt]{olplainarticle}
% Use option lineno for line numbers 

\title{Example Article Title}

\author[1]{Jakob Reinhardt}
\affil[1]{Yale University}



\begin{document}

\flushbottom
\thispagestyle{empty}

\section*{Research Question} How much should property be taxed? In a standard public finance model, goods that (a) have inelastic demands and (b) are disproportionately consumed by the wealthy should be taxed relatively more \citep{DiamondMirrlees1971a, DiamondMirrlees1971b}. As an implication, the large transaction costs involved in buying property and homeownership being more concentrated among the wealthy could suggest property should be taxed quite heavily. The goal of this project is to test this hypothesis and empirically estimate behavioral responses to property taxation. In a second step, outside of this class, I would use these estimates as an input to a structural model of optimal property taxation.\\
As property taxes in the US amount to more than 10 percent of all tax revenue collected, designing the property tax scheme is highly important. They are a very salient tax, they affect a large share of the income distribution, and are among the most controversial taxes (in the US). Hence, any insight on how they should be designed could be highly relevant to public discourse. Departing from \cite{Tiebout1956}, economists tend to frame property taxation as a price for the provision of local public goods, such as schools and infrastructure. Hence, an optimal tax trades off costs and benefits of providing public goods.  I want to use the behavioral responses that our current property tax system generates to make statements about how such a system should be designed. (As an aside, I am also considering framing this around property taxes as part of redistribution and not public goods provision - not sure which perspective is more insightful).\\
Several papers have addressed optimal property taxation. First, \cite{coven_etal} calibrates a rich overlapping generation model with financial constraints, migration and labor supply decisions, as well as rental markets to evaluate the effect of raising property taxes in California to the level of taxes in Texas. They find large positive effects on homeownership among younger generations, suggesting property taxes could be an important tool in tackling what is called the ``affordability crisis''. Second, \cite{balke_etal} calibrate a life-cycle model and reach broadly similar conclusions. Third, \cite{GieseckeMateen2022} analyze behavioral responses to property taxation, utilizing the large shock to home values in the aftermath of the great recession. They show local governments sharply increased the tax rate to keep budgets balanced, which caused some households to migrate. Broadly speaking, it seems the literature on optimal property taxation relies heavily on structural models, perhaps to address difficulties in finding estimates of behavioral responses to property taxation. By contrast, in this project I want to frame optimal tax tradeoffs in a sufficient statistics tradition, given that I can credibly identify behavioral responses.\\

\section*{Institutional Setting}
Key to this project is credibly identifying behavioral responses to property taxation. Consequently, I want to leverage two recent large tax reforms in the US. Prior to the reforms, State and local taxes (SALT) were deductible from the federal income tax base. In 2017, this option was capped at \$10,000, a limit that was later increased to \$40,000, roughly speaking\footnote{Starting at a modified adjusted gross income (MAGI) of \$500,000, the cap is phased down towards \$10,000. Hence, to be accurate, households with MAGI less than \$500,000 face a \$40,000 cap, households above \$600,000 MAGI face a cap of \$10,000, and other households face something in between.}. The deductibility limit introduces a sharp discontinuity in marginal property tax rates that households face, which increase by 20-50\% (depending on taxable income) at the threshold. This is a good setting for two reasons. First, the total amount of tax collected by local government remained unaffected, as the deduction can be seen as a change in transfers from the federal government. It implies the tax affected only the cost-side of the optimal tax tradeoff, not the benefits. This is key to separating the underlying elasticity of the demand for housing from behavioral responses to adjusting just the level of public goods. Second, there is large variation in tax rates across local governments. Hence, the cutoff of \$10,000 is located at very different parts of the income distribution across space, allowing to uncover a very broad range of behavioral responses as a function of income. For optimal tax purposes, this is very helpful, as the covariance between behavioral responses and income rank are important. Depending on the time horizon of this project, the second tax change in 2025 could then be leveraged to cross-check estimates across the two tax reforms for external validity.\\
There are other settings which might also reveal behavioral responses to property taxation. One setting I am exploring is variation in assessment levels in Germany. Assessment levels are indexed to assessed values in 1964; nowadays there is little correlation between property values in 1964 and 2025. This leads to quasi-random variation in effective property tax rates that households face. However, the setting being in Germany, it is very hard to access linked data on assessment values and household location decisions.\\

\section*{Data and Feasibility}
An ideal data set combines household-level data on itemization behavior from federal tax returns with data on households' locations and property ownership. The latter might be attainable using Yale's access to CoreLogic housing data and either data from DataAxle or Infutor. By contrast, based on what I heard from faculty, getting access to IRS tax return data seems to be very challenging in the current policy environment. One alternative is to reach out to a state-level department of revenue and work with a single, populous state. This would trade off less measurement error (itemization behavior is directly observed) with reductions in power and a smaller range of the income distribution covered. Another is to use aggregated data on itemization on the zip-code level provided by the IRS Statistics of Income (SOI) and assume homogeneous itemization rates conditional on zip code and taxable income. As a third option, I could rely only on aggregate data from the SOI and run regressions on the zip-code level, as previous papers have done. The drawback is that I don't actually see the distribution of effective marginal property tax rates.\\
Currently, the approach mixing microdata on property and location decisions with aggregate data on itemizations seems to be most promising to me. I could then use aggregate SOI data not used in estimation (for example the totals deducted) to cross-check my identification assumptions. Best that I can tell, all of this data (CoreLogic and DataAxle/Infutor) is accessible on Yale's high performance computing clusters and is mergable without IRB approval.\\

\section*{Research Design}
An ideal experiment would randomly assign different households different property tax schedules. Differences in marginal rates at a specific assessed value reveal intensive margin responses; differences in average tax rates reveal extensive margin responses in the form of cross-jurisdiction migration. I attempt to replicate such an experiment using the discontinuity in effective marginal property tax rates for households paying exactly \$10,000. Household above the threshold face an effective marginal tax rate up to 50\% higher than households located below. Comparing differences between these groups across time after the tax change in 2017 should then elicit the behavioral response I want to analyze.\\
The most important limitations are measurement error in taxable income and itemization behavior being unobserved. Measurement error in income implies I cannot set up and RDD at the deductibility threshold, as some observations will be on the wrong side of the cutoff. However, at any given income level, I can still accurately estimate the \textit{share} of people above/below the threshold, assuming mean-zero, homoskedastic measurement error. The second limitation is not observing itemization behavior. A household is only affected by the reform if they itemize deductions, and the standard deduction was increased from \$12,700 to \$24,000 for married couples filing jointly (which is the majority of households affected). To address this issue, I plan to use aggregated data on itemization behavior in place of individual-level data. Combining the share of people above the SALT deduction threshold and the share of itemizers, each conditional on income, will then identify the increase in average increase in effective property tax rates conditional on income, under some assumptions that I will spell out in more detail.\\ 
For estimation, I want to compare behavioral responses (denote the outcome of interest $Y_{it})$) across people who had different reform-induced changes in effective marginal property tax rates (denoted $\Delta \tau_{it}$ over time. This requires parallel trends between units differentially exposed to treatment, an assumption which has recently received much scrutiny \citep{Roth2023, GhanemSantannaWuthrich}. For now, I plan to address these concerns using the set-identification strategies laid out in \cite{Roth2023} and will think if I can come up with an explicit selection mechanism as motivated by \cite{GhanemSantannaWuthrich}. In any case, comparing trends in $Y_{it}$ conditional on $\Delta\tau_{it}$ will be very helpful in assessing the validity of parallel-trends based identification.\\
The main challenge is $\Delta \tau_{it}$ not being directly observed, because I lack individual-level data on itemizations. I want to circumvent this using \textit{aggregated} data on itemization, as the share of itemizers conditional on income and zip code is published in the IRS Statistics of Income (SOI). For identification this requires property taxes being independent on itemizations conditional on zip code and income\footnote{Towards a formal statement, suppose we observe data on income $y_{it}$, state income taxes $s_{it}$, total property taxes $p_{it}$, property tax rates $\tau_{it}$ and zip codes $z_{i}$ as well as aggregated data on itemizations $D_{it}$ conditional on zip code and taxable income, meaning $E\left[ D_{it} \mid z_i, y_{it}\right]$. We want to identify $E\left[\tau_{it} | z_i, y_{it}\right]$. Keep in mind that the individual effective property tax rate is the (observed) statutory rate modified by the (observed) marginal federal income tax rate $T\prime (y_{it})$ if a household itemizes (unobserved) and is not affected by the cap (observed). Also, $s_{it}$ is a known function of $y_{it}$. Adding a jurisdiction index $j$, we can then write $$E\left[\tau_{ijt} | z_i, y_{ijt}\right] = E\left[\tau_{jt} \mid z_i, y_{it}\right] - E\left[\tau_{ijt}D_{ijt}T\prime(y_{ijt})\mathbb{I}(s_{ijt} + p_{ijt} \leq 10,000) \mid z_i, y_{ijt}\right]$$ In words, average property tax rates in zip code $z_i$ at income $y_{ijt}$ are  the average nominal tax rate faced by residents with incomes $y_i$, modified by average marginal net of federal income tax rate for those who itemize and are not affected by the cap. The first term, average property taxes in a zip code conditional on income, is observed. For the second, the marginal federal income tax rate is observed, too and can be pulled outside of the equation, as can the local property tax rate. Identifying $E\left[D_{ijt}\mathbb{I}(s_{ijt} + p_{ijt} \leq 10,000)\right \mid z_i, y_{ijt}]$ from just $E\left[ D_{it} \mid z_i, y_{it}\right]$ is the remaining challenge, and follows if $D_{ijt} \perp \!\!\! \perp p_{ijt} \mid z_i, y_{ijt}$ holds.}. This is a questionable assumption, since households with higher property taxes due also have higher expenditures that could be itemized, thereby pushing them closer to the threshold of itemization reducing the federal tax burden. I have two ideas to address this. One is to compare observed itemized aggregates from the SOI to the ones implied by independence, the other is to apply a bounding exercise\footnote{This section is intentionally vague, as I have not thought through the details. The point is that the SOI contains data on which items are deducted, giving a bound in how important property taxes are in shaping itemization decisions.}. Under these assumptions, I can recover the change in effective property tax rates caused by the TCJA on a granular level.\\
Aside from identification, measurement error in income is a second major concern. Since I can merge multiple sources of income data, I am optimistic I can apply some results on estimation with measurement error from Econometrics III in this context.\\
In terms of implementation, this estimation is done on Yale's HPC, as both Data Axle and the joint CoreLogic is already accessible in a SQL database.\\


\bibliography{sample}

\end{document}